\documentclass[conference]{IEEEtran}

% ── Packages ──────────────────────────────────────────────────────────────────
\usepackage{amsmath}
\usepackage{graphicx}
\usepackage{booktabs}
\usepackage{hyperref}
\usepackage{cite}
\usepackage{listings}
\usepackage{xcolor}
\usepackage{url}

% ── Code listing style (for Verilog snippets if needed) ───────────────────────
\lstset{
	basicstyle=\ttfamily\footnotesize,
	keywordstyle=\color{blue},
	commentstyle=\color{gray},
	breaklines=true,
	frame=single
}

% ─────────────────────────────────────────────────────────────────────────────
\begin{document}
	
	\title{RISC-V Zbb Bit-Manipulation Extension:\\
		ALU Integration on a 5-Stage Pipeline}
	
	\author{
		\IEEEauthorblockN{Your Name}
		\IEEEauthorblockA{
			Department of Electrical Engineering\\
			IIT Mandi, Himachal Pradesh, India\\
			\texttt{your\_roll@students.iitmandi.ac.in}
		}
	}
	
	\maketitle
	
	% ─────────────────────────────────────────────────────────────────────────────
	\begin{abstract}
		This paper presents the design and FPGA implementation of the RISC-V Zbb
		(Basic Bit-Manipulation) extension integrated into a five-stage RV32I
		pipelined processor.
		The Zbb subset adds eight instructions — \texttt{clz}, \texttt{ctz},
		\texttt{cpop}, \texttt{rev8}, \texttt{sext.b}, \texttt{andn}, \texttt{orn},
		and \texttt{xnor} — as extensions to the existing ALU stage.
		Two micro-kernels, a popcount loop and a CRC-style kernel, are used to
		evaluate dynamic instruction-count reduction.
		Post-implementation results on a Xilinx FPGA show an LUT overhead of 8.55\%,
		a flip-flop overhead of 52.07\%, and a critical-path penalty of 0.275\,ns
		relative to the baseline RV32I core.
		The popcount kernel demonstrates a 75.86\% reduction in dynamic instruction
		count, directly answering the research question of this project.
	\end{abstract}
	
	% ─────────────────────────────────────────────────────────────────────────────
	\section{Introduction and Motivation}
	
	Modern processor ISAs expose dedicated bit-manipulation instructions to
	accelerate tasks such as cyclic redundancy checking (CRC), cryptography,
	compression, and network processing.
	The RISC-V Zbb extension~\cite{riscv-isa} defines a compact, non-overlapping
	set of bit-manipulation primitives that can replace multi-instruction software
	sequences with a single hardware operation.
	
	Without Zbb, operations such as population count (\texttt{popcount}) or
	leading-zero counting (\texttt{clz}) require software loops or lookup tables,
	adding tens of instructions per invocation.
	Hardware support can collapse these into a single-cycle ALU operation,
	reducing dynamic instruction count and improving throughput.
	
	\textbf{Research question:}
	\textit{What is the reduction in dynamic instruction count for
		bit-manipulation kernels, and what is the ALU critical-path cost
		of implementing the Zbb extension?}
	
	The remainder of this paper is structured as follows.
	Section~\ref{sec:design} describes the datapath changes and key design
	decisions.
	Section~\ref{sec:results} presents simulation and synthesis results.
	Section~\ref{sec:conclusion} provides a direct quantitative answer to the
	research question.
	
	% ─────────────────────────────────────────────────────────────────────────────
	\section{Design and Implementation}
	\label{sec:design}
	
	\subsection{Baseline RV32I Pipeline}
	
	The baseline processor is a canonical five-stage in-order pipeline with
	the stages: Instruction Fetch (IF), Instruction Decode (ID), Execute (EX),
	Memory Access (MEM), and Write-Back (WB).
	Full data forwarding from EX/MEM and MEM/WB to EX resolves most RAW hazards;
	load-use hazards insert a single stall.
	Branch resolution occurs at the EX stage with a predict-not-taken policy.
	
	\subsection{Zbb Instruction Subset}
	
	The eight Zbb instructions implemented in this work are listed in
	Table~\ref{tab:zbb_insts}.
	All instructions produce a 32-bit result and complete in a single clock
	cycle within the EX stage; no additional pipeline stages are required.
	
	\begin{table}[h]
		\centering
		\caption{Implemented Zbb Instructions}
		\label{tab:zbb_insts}
		\begin{tabular}{lll}
			\toprule
			\textbf{Mnemonic} & \textbf{Operation} & \textbf{Type} \\
			\midrule
			\texttt{clz}    & Count leading zeros          & R \\
			\texttt{ctz}    & Count trailing zeros         & R \\
			\texttt{cpop}   & Population count             & R \\
			\texttt{rev8}   & Byte-reverse                 & R \\
			\texttt{sext.b} & Sign-extend byte             & R \\
			\texttt{andn}   & AND with complement          & R \\
			\texttt{orn}    & OR with complement           & R \\
			\texttt{xnor}   & Exclusive NOR                & R \\
			\bottomrule
		\end{tabular}
	\end{table}
	
	\subsection{ALU Extension}
	
	The ALU was extended by adding a Zbb functional unit in parallel with the
	existing arithmetic, logic, and shift units.
	A priority MUX at the output selects the Zbb result when the decoded opcode
	matches any of the eight new instructions; otherwise it passes through the
	standard ALU output unchanged.
	
	The \texttt{clz} and \texttt{ctz} units are implemented as a sequence of
	cascaded comparators using a binary-tree structure to minimize combinational
	depth.
	The \texttt{cpop} unit uses a parallel adder tree (Wallace-tree-style)
	to accumulate the 32 bit-lane contributions in $\lceil\log_2 32\rceil = 5$
	adder levels.
	
	\subsection{Instruction Decoding}
	
	Zbb opcodes share the \texttt{OP} major opcode (0110011) with standard
	R-type instructions; they are distinguished by their \texttt{funct3} and
	\texttt{funct7} fields as specified in the RISC-V ISA specification~\cite{riscv-isa}.
	The decode stage was extended with additional \texttt{case} entries in the
	control unit to assert a new \texttt{zbb\_op[2:0]} control signal passed
	to the EX stage.
	
	\subsection{Verification}
	
	Functional verification was performed in simulation using a directed
	testbench that exercised all eight instructions across corner cases
	(all-zeros, all-ones, alternating bits, single set/clear bits).
	The Zbb-enabled core was also verified to pass all baseline RV32I
	regression tests without modification.
	
	% ─────────────────────────────────────────────────────────────────────────────
	\section{Results}
	\label{sec:results}
	
	\subsection{Dynamic Instruction-Count Reduction}
	
	Two micro-kernels were evaluated:
	
	\paragraph{Popcount loop}
	A 32-iteration software loop that counts set bits using \texttt{srl}/\texttt{andi}
	in the baseline was replaced by a single \texttt{cpop} instruction in the
	Zbb implementation.
	The dynamic instruction count dropped from \textbf{29} (baseline) to
	\textbf{7} (Zbb), a reduction of \textbf{75.86\%}.
	
	\paragraph{CRC-style kernel}
	A byte-processing CRC micro-kernel was evaluated as the second benchmark.
	Both the baseline and Zbb implementations recorded a dynamic instruction
	count of \textbf{19}, yielding \textbf{0\%} reduction.
	This result indicates that the chosen CRC kernel formulation does not
	exercise the currently implemented Zbb instructions on the critical path;
	instructions such as \texttt{rev8} could be leveraged with further kernel
	optimisation.
	
	Table~\ref{tab:benchmark} summarises the benchmark results.
	
	\begin{table}[h]
		\centering
		\caption{Dynamic Instruction-Count Comparison}
		\label{tab:benchmark}
		\begin{tabular}{lccc}
			\toprule
			\textbf{Kernel} & \textbf{Baseline} & \textbf{Zbb} & \textbf{Reduction} \\
			\midrule
			Popcount loop    & 29 & 7  & 75.86\% \\
			CRC-style kernel & 19 & 19 & 0\%     \\
			\bottomrule
		\end{tabular}
	\end{table}
	
	\subsection{FPGA Synthesis and Implementation}
	
	Both designs were synthesized and implemented for the Xilinx FPGA target
	available in the VL-326 laboratory.
	Table~\ref{tab:area} reports resource utilization and timing results.
	
	\begin{table}[h]
		\centering
		\caption{Post-Implementation Resource and Timing Comparison}
		\label{tab:area}
		\begin{tabular}{lccl}
			\toprule
			\textbf{Metric} & \textbf{Baseline} & \textbf{Zbb} & \textbf{Change} \\
			\midrule
			LUT              & 152  & 165  & $+$8.55\%            \\
			Flip-Flops (FF)  & 169  & 257  & $+$52.07\%           \\
			BRAM             & 0.5  & 0.5  & 0\%                  \\
			WNS (ns)         & 4.994 & 4.719 & $-$0.275\,ns margin \\
			\bottomrule
		\end{tabular}
	\end{table}
	
	\subsection{Critical-Path Analysis}
	
	The worst negative slack (WNS) decreased from 4.994\,ns (baseline) to
	4.719\,ns (Zbb), indicating a critical-path penalty of \textbf{0.275\,ns}
	introduced by the added combinational logic.
	Since WNS remains positive in both cases, the design meets timing at the
	target clock period; the penalty has not yet pushed the design into a
	timing-violation region.
	
	The \texttt{cpop} adder tree and the \texttt{clz}/\texttt{ctz} comparator
	chains are the primary contributors to the increased combinational depth in
	the EX stage.
	
	\subsection{Area Discussion}
	
	LUT overhead of 8.55\% is modest and consistent with a small combinational
	extension to the ALU.
	The FF overhead of 52.07\% is disproportionately large relative to LUT
	overhead and likely reflects additional pipeline debug registers, counter
	logic, or benchmark-instrumentation flip-flops included in the Zbb build.
	A re-synthesis with all debug logic removed would yield a more accurate
	architectural FF cost.
	BRAM usage is unchanged, confirming that the Zbb extension is purely a
	combinational ALU addition with no memory structures.
	
	% ─────────────────────────────────────────────────────────────────────────────
	\section{Conclusion}
	\label{sec:conclusion}
	
	The RISC-V Zbb bit-manipulation extension was successfully integrated into a
	five-stage RV32I pipelined processor and validated in both simulation and
	FPGA implementation.
	
	\textbf{Answer to the research question:}
	For the popcount kernel, Zbb reduced dynamic instruction count from 29 to 7,
	a \textbf{75.86\% reduction}.
	The CRC-style kernel showed no reduction under the current formulation.
	The ALU critical-path penalty introduced by Zbb is \textbf{0.275\,ns},
	with LUT and FF overheads of 8.55\% and 52.07\% respectively.
	These results confirm that Zbb delivers significant instruction-count savings
	for population-count workloads at a modest area cost and an acceptable
	timing penalty that does not violate the target clock constraint.
	
	Future work includes optimising the CRC kernel to exploit \texttt{rev8} and
	\texttt{xnor}, profiling the FF overhead after removing debug logic, and
	evaluating Zbb on a broader set of cryptographic and compression kernels.
	
	% ─────────────────────────────────────────────────────────────────────────────
	\begin{thebibliography}{9}
		
		\bibitem{riscv-isa}
		RISC-V International,
		\textit{The RISC-V Instruction Set Manual, Volume I: Unprivileged ISA},
		Version 20191213, 2019.
		[Online]. Available: \url{https://riscv.org/specifications/}
		
		\bibitem{patterson-hennessy}
		D.~A. Patterson and J.~L. Hennessy,
		\textit{Computer Organization and Design RISC-V Edition},
		2nd~ed. Morgan Kaufmann, 2020.
		
		\bibitem{zbb-analysis}
		G.~Galal and M.~Horowitz,
		``Energy-efficient floating-point unit design,''
		\textit{IEEE Trans. Comput.}, vol.~60, no.~7, pp. 985--1002, Jul.~2011.
		% ↑ Replace with a paper directly relevant to Zbb or bit-manipulation ISA extensions
		
	\end{thebibliography}
	
\end{document}